% Transcription — fonds Alexandre Grothendieck, shelfmark 115, batch 1 (pages 1-14).
% Established from the facsimile archives/batches/115/batch-01.pdf (a mirror of
% https://grothendieck.umontpellier.fr/115.pdf), per the skill
% transcribe-grothendieck. Pages 4, 6 and 12 are versos of unrelated computer
% listings (dated 02 JUN 82) — absent, as they should be.
% Pass: Fable 5 (claude-fable-5), 2026-08-08 — first pass, unchecked against the pages by a human.
\documentclass[11pt,a4paper]{article}
% Le préambule commun à toutes les transcriptions du fonds Grothendieck.
%
% Il tient en une idée : **l'incertitude se compose, elle ne se résout pas.**
% Une transcription qui lisse un mot douteux a détruit la seule chose pour
% laquelle on la faisait — un lecteur doit pouvoir savoir, à chaque endroit,
% ce qui était sur le feuillet et ce qui a été ajouté. Les six macros qui
% suivent sont donc l'appareil critique, et non de la décoration.
%
% Les mêmes commandes sont reconnues par `scripts/render.mjs`, qui produit la
% vue de lecture du site. Une macro ajoutée ici doit l'être aussi là-bas,
% faute de quoi le rendu échouera bruyamment — ce qui est voulu : un rendu
% silencieusement faux est pire qu'un rendu absent.

\ProvidesPackage{grothendieck}[2026/08/08 Transcriptions du fonds Grothendieck]

\usepackage{amsmath,amssymb,amsthm}
\usepackage{mathrsfs}
\usepackage{stmaryrd}
% Les diagrammes commutatifs sont partout dans ce fonds : c'est la langue
% même de la géométrie algébrique de Grothendieck, et une transcription qui
% les rendrait en prose ne transcrirait rien.
\usepackage{tikz-cd}
% Français par défaut — c'est la langue des feuillets — mais l'anglais est
% chargé aussi : la traduction et le résumé sont des documents anglais, et la
% typographie française (espaces avant les deux-points) y serait une faute.
% Ces documents basculent par \selectlanguage{english} en tête de corps.
\usepackage[english,french]{babel}
\usepackage{fontspec}
\usepackage{geometry}
\geometry{a4paper,margin=25mm,marginparwidth=28mm}
\usepackage{marginnote}
\usepackage{xcolor}
% ulem, not soul. `soul`'s \st and \ul are built on a character-by-character
% scan that XeLaTeX's font machinery breaks: the document fails with an
% undefined control sequence rather than degrading. ulem does the same two jobs
% and is Unicode-engine safe. `normalem` stops it hijacking \emph, which the
% transcriptions use for Grothendieck's own emphasis.
\usepackage[normalem]{ulem}
\usepackage{hyperref}
\hypersetup{colorlinks=true,linkcolor=black,urlcolor=black,pdfborder={0 0 0}}

% --- Métadonnées du lot -----------------------------------------------------
% Reprises telles quelles par le script de rendu, qui les affiche en tête de
% la vue de lecture. Elles fixent ce que le fichier prétend couvrir : sans
% elles, une transcription flotte sans point d'attache dans un fonds de seize
% mille feuillets.

\newcommand{\folder}[1]{\gdef\grothFolder{#1}}
\newcommand{\batch}[1]{\gdef\grothBatch{#1}}
\newcommand{\leaves}[2]{\gdef\grothFirst{#1}\gdef\grothLast{#2}}
\newcommand{\foldertitle}[1]{\gdef\grothFoldertitle{#1}}
\gdef\grothFolder{?}\gdef\grothBatch{?}\gdef\grothFirst{?}\gdef\grothLast{?}\gdef\grothFoldertitle{}

% --- Le résumé --------------------------------------------------------------
%
% Il ouvre la lecture modernisée, sous le titre « Résumé », et il est composé
% en petit. Ce n'est pas un ornement : c'est le seul passage du document qui
% ne soit pas des mathématiques, et un lecteur doit pouvoir le reconnaître
% avant de l'avoir lu — puis le sauter s'il connaît déjà le sujet, ou s'y
% arrêter s'il ne le connaît pas. Le filet qui le suit marque la reprise du
% corps.

\newenvironment{resume}
  {\small\setlength{\parskip}{0.45em}}
  {\par\medskip\hrule\medskip}

% --- Mots-clés ---------------------------------------------------------------
%
% En anglais, en clôture du résumé de la lecture modernisée : le vocabulaire
% moderne sous lequel chercher ce que les feuillets construisent. Le manifeste
% du site (`npm run manifest`) les extrait pour étiqueter la cote entière —
% cette ligne est la source unique des tags, il n'y a pas de fichier à côté.
\newcommand{\keywords}[1]{\par\smallskip\noindent
  {\footnotesize\textsc{Keywords} --- \textit{#1}}\par}

% --- L'appareil critique ----------------------------------------------------

% Le numéro de feuillet, en marge. C'est l'ancre qui permet de régler un
% désaccord sur une formule en regardant *un* feuillet plutôt qu'une chemise
% de six cents. Le site s'en sert aussi pour tourner les pages du fac-similé
% quand on fait défiler la transcription.
\newcommand{\leaf}[1]{%
  \marginnote{\footnotesize\color{gray!70}\textsf{#1}}%
  \label{leaf:#1}\ignorespaces}

% Illisible : on ne devine pas. Un mot inventé qui se lit comme les autres est
% le pire résultat possible.
\newcommand{\ill}{\textcolor{red!60!black}{[\,\ldots\,]}}

% Lecture douteuse : le mot est donné, et signalé comme incertain.
\newcommand{\uncertain}[1]{\uline{#1}}

% Ajout de l'éditeur — une lettre manquante, un mot sous-entendu.
\newcommand{\add}[1]{\textcolor{blue!50!black}{[#1]}}

% Biffé par Grothendieck lui-même. Ce qu'il a rayé fait partie du document :
% une rature dit souvent où la pensée a changé de direction.
\newcommand{\struck}[1]{\sout{#1}}

% Note du transcripteur, sur l'état matériel du feuillet.
\newcommand{\note}[1]{\par\smallskip{\small\color{gray!60!black}\textit{[#1]}}\par\smallskip}

% Note marginale de Grothendieck, distincte du corps du texte.
\newcommand{\marginal}[1]{\par\smallskip\noindent\hspace{1em}%
  {\small\color{gray!60!black}#1}\par\smallskip}

% --- Titre ------------------------------------------------------------------

\AtBeginDocument{%
  \begin{center}
    {\large\bfseries Fonds Alexandre Grothendieck --- cote n\textsuperscript{o}~\grothFolder}\\[2pt]
    {\itshape\grothFoldertitle}\\[4pt]
    {\small feuillets \grothFirst--\grothLast\ (lot \grothBatch)}\\[2pt]
    {\footnotesize\color{gray!60!black}Transcription établie d'après le fac-similé numérisé
     par l'Université de Montpellier.\\
     Les lectures incertaines sont soulignées, les illisibles notées [\,\ldots\,],
     les ajouts entre crochets.}
  \end{center}
  \vspace{1em}}

\endinput


\folder{115}
\batch{1}
\pages{1}{14}
\dating{[à partir de 1982]}
\foldertitle{« Correspondances » fonctorielles. Dualité des topos : notes manuscrites (s.d.)}

\begin{document}

\section*{Correspondances fonctorielles (pages 1 à 5)}

\page{1}
Titre porté à l'encre sur la chemise : « Correspondances » fonctorielles —
Dualité des topos. \note{chemise du dossier}

\page{2}
Formules d'adjonction, en tête de page :
\[
H_{A}\bigl(f^{*}(\Psi_{\varepsilon}),\ \Phi^{\varepsilon}\bigr)
\;\simeq\;
H_{B}\bigl(\Psi_{\varepsilon},\ f^{\varepsilon}_{!}(\Phi^{\varepsilon})\bigr)
\]
\note{lecture incertaine des décorations $\varepsilon$ et de la répartition
des indices}
\[
H_{A}\bigl(f^{*}(b),\ a\bigr) \;\simeq\; H_{B}\bigl(b,\ f^{\varepsilon}(a)\bigr)
= \mathrm{Hom}_{B}\bigl(b,\ f(a)\bigr),
\qquad
f^{\varepsilon}(a) = \varinjlim \cdots
\]
\note{le système d'indices de la colimite est illisible. Le reste de la page
— un cercle de catégories de foncteurs et de petits diagrammes — est un
premier état raturé de la roue de la page 5 ; lectures trop incertaines pour
être transcrites}

\page{3}
\note{page 1 de la main de l'auteur}
En tête, à droite : $\underline{\mathrm{Hom}}_{!}$ : sous-catégorie pleine de
$\underline{\mathrm{Hom}}$ formée des foncteurs commutant aux
\uncertain{limites inductives} ; $\underline{\mathrm{Hom}}^{!}$ :
sous-catégorie pleine de $\underline{\mathrm{Hom}}$ formée des foncteurs
commutant aux \uncertain{limites projectives}.

En marge gauche : cas $B = A^{\circ}$ : d'où l'élément canonique
$H = H_{A} \in (A \times A^{\circ})^{\wedge}$ ; \uncertain{les prolongements
canoniques} $\widehat{A}^{\circ} \hookrightarrow A^{\vee}$, \ldots
\note{le reste de la colonne marginale est trop raturé pour être transcrit}

Le diagramme central, simplifié — l'original est un espalier d'une dizaine de
flèches marquées $\approx$, dont plusieurs biffées :

\begin{tikzcd}[column sep=small, row sep=small, nodes={font=\scriptsize}]
\underline{\mathrm{Hom}}(B^{\circ}, \widehat{A})
  & (A\times B)^{\wedge} \arrow[l, "\approx"'] \arrow[r, "\approx"] \arrow[d, "\approx"]
  & \underline{\mathrm{Hom}}(A^{\circ}, \widehat{B}) \arrow[d, "\approx"] \\
  & \underline{\mathrm{Hom}}^{!!}\bigl(\widehat{A}^{\circ} \times \widehat{B}^{\circ},\ (\mathrm{Ens})\bigr) = \mathrm{Bifais}(\widehat{A}, \widehat{B})
  & \underline{\mathrm{Hom}}^{!}(\widehat{A}^{\circ}, \widehat{B}) \\
  & \underline{\mathrm{Hom}}_{!}(A^{\vee}, \widehat{B}) \arrow[u, "\approx"']
  &
\end{tikzcd}

Le dispositif systématique : $A \times B$, avec les projections
$p_{B} : A\times B \to A$ et $p_{A} : A\times B \to B$.
\note{l'indice de la projection nomme le facteur supprimé — c'est la seule
lecture sous laquelle les formules qui suivent concordent entre elles}
\[
\underline{\mathrm{Hom}}_{A}(X,Y) \overset{\mathrm{def}}{=}
   p_{B*}\,\underline{\mathrm{Hom}}_{A\times B}(X,Y),
\qquad
\underline{\mathrm{Hom}}_{B}(X,Y) \overset{\mathrm{def}}{=}
   p_{A*}\,\underline{\mathrm{Hom}}_{A\times B}(X,Y)
\]
\[
\mathrm{Hom}_{A\times B}(X,Y)
\simeq \Gamma_{A\times B}\,\underline{\mathrm{Hom}}_{A\times B}(X,Y)
\simeq \Gamma_{A}\,\underline{\mathrm{Hom}}_{A}(X,Y)
\simeq \Gamma_{B}\,\underline{\mathrm{Hom}}_{B}(X,Y),
\qquad X, Y \in (A\times B)^{\wedge}
\]
\[
{}^{s}X(L) = \underline{\mathrm{Hom}}_{B}\bigl(p_{B}^{*}L,\ X\bigr)
           = p_{A*}\,\underline{\mathrm{Hom}}\bigl(p_{B}^{*}L,\ X\bigr),
\qquad
{}^{d}X(M) = \underline{\mathrm{Hom}}_{A}\bigl(p_{A}^{*}M,\ X\bigr)
           = p_{B*}\,\underline{\mathrm{Hom}}\bigl(p_{A}^{*}M,\ X\bigr)
\]
(en marge : $L \in \widehat{A}$, $M \in \widehat{B}$)
\[
X(L,M) = \mathrm{Hom}\bigl(L \boxtimes M,\ X\bigr)
\simeq \mathrm{Hom}\bigl(L,\ {}^{d}X(M)\bigr)
\simeq \mathrm{Hom}\bigl(M,\ {}^{s}X(L)\bigr),
\qquad\text{où } L \boxtimes M = p_{B}^{*}L \times p_{A}^{*}M .
\]

« Ces formules sont particulièrement intéressantes pour
$L = a \in \mathrm{Ob}\,A$, $M = b \in \mathrm{Ob}\,B$. »
\note{lecture incertaine des deux membres de droite}

« Soit $X$ ; définissons $B^{\circ} \to \widehat{A}$ (savoir
$b \mapsto {}^{d}X(b) = \underline{\mathrm{Hom}}_{A}(p_{A}^{*}(b), X)$)
\ldots{} on trouve $B^{\circ} \to \mathrm{Ens}$ : c'est \ill{}, qui
s'écrit aussi $\underline{\mathrm{Hom}}_{B}(p_{B}^{*}(L), X) = {}^{s}X(L)$.
Interprétation \uncertain{duale} pour ${}^{d}X(M)$. »
\note{passage interligné, lecture partiellement conjecturale}

« $(L,M) \mapsto L \boxtimes M : \widehat{A} \times \widehat{B} \to
(A\times B)^{\wedge}$ est pleinement fidèle. »

« le \ill{} qui est le foncteur canonique \ldots{} dit « de Yoneda » pour
$\widehat{A} \times \widehat{B}$ »
\[
\underline{\mathrm{Hom}}_{A\times B}\bigl(L\boxtimes M',\ L'\boxtimes N'\bigr)
\simeq \underline{\mathrm{Hom}}_{A}(L,L') \boxtimes
       \underline{\mathrm{Hom}}_{B}(M',N')
\qquad \text{p.ex.}
\]
\[
\underline{\mathrm{Hom}}_{A\times B}\bigl(L\boxtimes M',\ L'\boxtimes e_{B}\bigr)
\simeq \underline{\mathrm{Hom}}_{A}(L,L') \boxtimes e_{B}
= p_{B}^{*}\,\underline{\mathrm{Hom}}_{A}(L,L')
\]
en particulier, faisant $L' = e_{A}$ : \ldots
\note{la suite de la ligne est illisible}

$\underline{\mathrm{Hom}}_{A\times B}\bigl(p_{A}^{*}(M), \cdots\bigr)
\Leftarrow p_{B}^{*}\cdots$ (les flèches étant les flèches canoniques)
\note{ligne partiellement illisible}

Posant $X^{M} \overset{\mathrm{def}}{=}
\underline{\mathrm{Hom}}_{A\times B}\bigl(p_{A}^{*}M,\ X\bigr)$, on trouve
donc $(p_{B}^{*}L)^{M} \simeq p_{B}^{*}L\cdots$
\note{fin de formule incertaine}

\page{5}
\note{page 2 de la main de l'auteur}
En marge haute, un recensement : $A^{\circ}, \widehat{A}^{\circ}, A^{\vee}$
et $B^{\circ}, \widehat{B}^{\circ}, B^{\vee}$, reliés par des flèches
croisées à $\widehat{A}$, $\widehat{B}$, avec la mention « (6 cas) » ;
au-dessous,
$\bigl|\widehat{A}^{\circ} \text{ ou } A^{\vee}\bigr| \times
 \bigl|\widehat{B}^{\circ} \text{ ou } B^{\vee}\bigr| \to \mathrm{Ens}$,
avec la mention « (4 cas) ».
\note{lu d'abord « (4 co) », « (6 (co)) », non élucidé ; la lecture est
« (4 cas) », « (6 cas) », et les deux mentions dénombrent la roue
ci-dessous : quatre formes à trois arguments, six formes à deux arguments}

\emph{NB} en marge : $A^{\circ} \to \widehat{A}^{\circ}$,
$A^{\circ} \to A^{\vee}$ ; $B^{\circ} \to \widehat{B}^{\circ}$,
$B^{\circ} \to B^{\vee}$.

La roue des catégories de foncteurs autour de
$(A\times B)^{\wedge} = \underline{\mathrm{Hom}}(A^{\circ}, B^{\circ};\,
\mathrm{Ens})$ — un cercle porté par quatre rayons, avec une couronne
intérieure et un bord :

\begin{tikzcd}[column sep=tiny, row sep=small, nodes={font=\scriptsize}]
(\widehat{B}_{\to}) & & & & & & (\widehat{A}_{\to}) \\
& \underline{\mathrm{Hom}}^{!}(\widehat{B}^{\circ},\widehat{A}) \arrow[rr, bend left=22] \arrow[dd, bend right=22] & & \underline{\mathrm{Hom}}^{!!}(\widehat{A}^{\circ},\widehat{B}^{\circ};\,\mathrm{Ens}) & & \underline{\mathrm{Hom}}^{!}(\widehat{A}^{\circ},\widehat{B}) \arrow[ll, bend right=22] \arrow[dd, bend left=22] \arrow[dl] & \\
& & \underline{\mathrm{Hom}}(B^{\circ},\widehat{A}) \arrow[ul] & & \underline{\mathrm{Hom}}(A^{\circ},\widehat{B}) & & \\
& \underline{\mathrm{Hom}}_{!}^{!}(A^{\vee},\widehat{B}^{\circ};\,\mathrm{Ens}) & & \substack{(A\times B)^{\wedge} \\ =\, \underline{\mathrm{Hom}}(A^{\circ},B^{\circ};\,\mathrm{Ens})} \arrow[ul] \arrow[ur] \arrow[dl] \arrow[dr] \arrow[uu, no head, "\wedge\wedge" description] \arrow[ll, no head, "\vee\wedge" description] \arrow[rr, no head, "\wedge\vee" description] \arrow[dd, no head, "\vee\vee" description] & & \underline{\mathrm{Hom}}^{!}_{!}(\widehat{A}^{\circ},B^{\vee};\,\mathrm{Ens}) & \\
& & \underline{\mathrm{Hom}}(A^{\circ},\widehat{B}) \arrow[dl] & & \underline{\mathrm{Hom}}(B^{\circ},\widehat{A}) \arrow[dr] & & \\
& \underline{\mathrm{Hom}}_{!}(A^{\vee},\widehat{B}) \arrow[uu, bend left=22] \arrow[rr, bend right=22] & & \underline{\mathrm{Hom}}_{!!}(A^{\vee},B^{\vee};\,\mathrm{Ens}) & & \underline{\mathrm{Hom}}_{!}(B^{\vee},\widehat{A}) \arrow[uu, bend right=22] \arrow[ll, bend left=22] & \\
(A^{\vee}_{\to}) & & & & & & (B^{\vee}_{\to})
\end{tikzcd}

« toutes les flèches du diagramme \ldots{} sont des équivalences de
catégories »
\note{la roue de la page, relue sur le fac-similé à haute résolution : les
quatre rayons (axes vertical et horizontal) sont des traits fins sans
pointe, portant les marques de variance $\wedge\wedge$, $\vee\wedge$,
$\wedge\vee$, $\vee\vee$ à mi-course ; les quatre annotations d'angle
$(\widehat{B}_{\to})$, $(\widehat{A}_{\to})$, $(A^{\vee}_{\to})$,
$(B^{\vee}_{\to})$ sont hors du cercle, une par quadrant. Onze catégories
distinctes : le centre ; six formes à deux arguments (les « 6 cas » de la
marge — les deux formes intérieures sont chacune écrites deux fois, aux
deux bouts d'un diamètre) ; quatre formes mixtes à trois arguments (les
« 4 cas »), aux quatre points cardinaux. Les huit arcs du bord, appuyés,
convergent tous vers ces quatre formes mixtes. Trois diagonales sortent du
centre vers le bord ; dans le seul quadrant supérieur droit la flèche fine
entre bord et couronne rentre,
$\underline{\mathrm{Hom}}^{!}(\widehat{A}^{\circ},\widehat{B}) \to
\underline{\mathrm{Hom}}(A^{\circ},\widehat{B})$ — une restriction là où
les trois autres sont dessinées comme des prolongements. Une première
lecture avait réduit la roue à huit flèches rayonnantes ; le dessin
ci-dessus suit désormais la page trait par trait}

En bas — familles d'indices : « $I$ ens.\ d'indices, $(A_i)_{i\in I}$ famille
de (petites) catégories. Si $I' \subset I$, on pose
$A_{I'} = \prod_{i\in I'} A_i$. »
\[
I = \coprod_{j\in J} I_{j}, \qquad J = J^{+} \amalg J^{-},
\qquad I^{+} = \coprod_{j\in J^{+}} I_{j}, \quad
I^{-} = \coprod_{j\in J^{-}} I_{j}
\]
« \uncertain{Je suppose} $I$, $J$ finis pour « plus safe ». »
\[
\widehat{A_{I}} \;\simeq\;
\underline{\mathrm{Hom}}^{!!}\Bigl(\prod_{j\in J^{-}} A_{I_{j}}^{\vee}\,,\
\prod_{j'\in J^{+}} \widehat{A_{I_{j'}}}^{\circ}\,;\ \mathrm{Ens}\Bigr)
\]
avec des flèches $\approx$ vers
\[
\underline{\mathrm{Hom}}_{!}\Bigl(\prod_{j\in J^{-}} A_{I_{j}}^{\vee},\
\widehat{A_{I^{+}}}\Bigr),
\qquad
\underline{\mathrm{Hom}}^{!}\Bigl(\prod_{j'\in J^{+}}
\widehat{A_{I_{j'}}}^{\circ},\ A_{I^{-}}^{\vee}\Bigr)
\]
\note{variances partiellement conjecturales}

\section*{La dualité entre $\widehat{A}$ et $A^{\vee}$ (pages 7 à 9)}

\page{7}
\note{page 3 de la main de l'auteur}
\[
\boxed{\ H_{A}(\Phi,\Psi)
= \mathrm{Hom}_{\widehat{A}}\bigl(\Phi,\ \Psi_{A}(\Psi)\bigr)
\simeq \mathrm{Hom}_{A^{\vee}}\bigl(\Psi,\ \varphi_{A}(\Phi)\bigr)\ }
\]
(en marge : $\widehat{A}^{\circ} \xrightarrow{\varphi_{A}} A^{\vee}$,
$A^{\vee\circ} \xrightarrow{\Psi_{A}} \widehat{A}$,
$\widehat{A} \times A^{\vee\circ} \xrightarrow{H_{A}} \mathrm{Ens}$)

« i.e.\ on a $A^{\vee\circ} \rightleftarrows \widehat{A}$
($\varphi_{A}^{\circ}$, $\Psi_{A}$), paire de foncteurs adjoints. »

NB — a) le triangle
\begin{tikzcd}[column sep=small, row sep=small]
A^{\vee\circ} \arrow[rr, "\Psi_{A}"] & & \widehat{A} \\
& A \arrow[ul, hook] \arrow[ur, hook'] &
\end{tikzcd}
est commutatif, et $\Psi_{A}$ commute aux \uncertain{limites projectives} —
ces \uncertain{deux} conditions \uncertain{caractérisent} $\Psi_{A}$ à iso.
unique près ; — b) de même pour
$\widehat{A} \xrightarrow{\varphi_{A}^{\circ}} A^{\vee\circ}$ :
« commutatif, et $\varphi_{A}^{\circ}$ commute aux \ill{} et \ill{} — et ces
conditions : $\varphi_{A}^{\circ}$ à iso.\ unique près. »

En marge gauche : \uncertain{$\varphi_{A^{\circ}} = \Psi_{A}$},
\uncertain{$\Psi_{A^{\circ}} = \varphi_{A}$},
\uncertain{$H_{A^{\circ}}(\Psi,\Phi) = H_{A}(\Phi,\Psi)$}
(\uncertain{formules à vérifier}).

Morphismes d'adjonction :
\[
\Phi \longrightarrow \Psi_{A}\,\varphi_{A}^{\circ}(\Phi)
\ \text{dans } \widehat{A},
\qquad
\Psi \longrightarrow \varphi_{A}\,\Psi_{A}^{\circ}(\Psi)
\ \text{dans } A^{\vee} ;
\]
« iso si $\Phi \in \mathrm{Ob}\,A$ \uncertain{ou} $\Psi \in \mathrm{Ob}\,A$ ».

c)
\[
\boxed{\ H_{A}(a,b) = \mathrm{Hom}_{A}(a,b)\ }
\]
« donc $H_{A}$ prolonge $\mathrm{Hom}_{A}$, et est \uncertain{déterminé} à
iso.\ unique près par cette propriété, plus celle que $H_{A}$ commute aux
\uncertain{limites} \ldots »
\[
H_{A}(\Phi, b) \simeq \mathrm{Hom}_{\widehat{A}}(\Phi,\ b),
\qquad
H_{A}(a, \Psi) \simeq \mathrm{Hom}_{A^{\vee}}(\Psi,\ a)
= \mathrm{Hom}_{A^{\vee\circ}}(a,\ \Psi)
\]
\note{seconde formule partiellement conjecturale}

« Question : que penser des morphismes d'adjonction $\alpha_{\Phi}$,
$\beta_{\Psi}$ ? Iso \uncertain{(oui)} seulement si $\Phi$ \uncertain{resp.}
$\Psi$ sont représentables ? Est-ce \uncertain{que c'est} une équivalence
faible (sans doute pas). Les foncteurs $\varphi_{A}$ et $\Psi_{A}$ sont-ils
\ill{} ? »

« Relations avec \uncertain{les contractions} : »
\[
\boxed{\ H^{A} :\ (\Phi,\Psi) \longmapsto \Phi \star_{A} \Psi\ ;\quad
\widehat{A} \times A^{\vee} \longrightarrow \mathrm{Ens}\ }
\]
d)
\[
\boxed{\ H^{A}(a,b) = \mathrm{Hom}_{A}(b,a)\ }
\]
« et $H^{A}$ commute aux \uncertain{limites inductives}, en \uncertain{quelles}
conditions : iso.\ unique près. »
\[
\boxed{\ \widehat{A} \xrightarrow{\ \sim\ }
\underline{\mathrm{Hom}}_{!}(A^{\vee}, \mathrm{Ens}),
\qquad
A^{\vee} \xrightarrow{\ \sim\ }
\underline{\mathrm{Hom}}_{!}(\widehat{A}, \mathrm{Ens})\ }
\]
\note{la décoration du second Hom diffère du premier — point d'exclamation
placé autrement, lecture incertaine}
\[
\Phi \star a = \Phi(a), \qquad a \star \Psi = \Psi(a)
\]
\[
H^{A}(\Phi,\Psi) = \varphi^{A}(\Phi)(\Psi) \simeq \Psi^{A}(\Psi)(\Phi)
\]
\note{le premier membre est précédé d'un symbole non élucidé}

\page{8}
\note{page entière biffé de trois grands traits obliques}
Multivariances :
\[
\Bigl(\prod_{i\in I} A_{i}\Bigr)^{\wedge},
\qquad
A_{1}^{\circ} \times A_{2}^{\circ} \times \cdots \times A_{n}^{\circ}
\to \mathrm{Ens},
\qquad
\widehat{A_{1}}^{\circ} \times \widehat{A_{2}} \times A_{3}^{\vee}
\times \cdots
\]
\note{la troisième liste de variances est partiellement illisible}
\[
X_{1} \otimes X_{2} \simeq \mathrm{Hom}(X_{1}^{\vee}, X_{2})
\simeq \mathrm{Hom}(X_{2}^{\vee}, X_{1})
\]
\note{un membre intermédiaire, portant un 1 souligné, reste incertain}
\[
F(X_{1}\boxtimes X_{2})
\simeq \underline{\mathrm{Hom}}^{!!}\bigl(F(X_{1})^{\circ},\ F(X_{2})^{\circ};\
\mathrm{Ens}\bigr)
\simeq \cdots
\simeq \mathrm{Hom}\bigl(F(X_{1})^{\vee},\ F(X_{2})\bigr)
\simeq \mathrm{Hom}\bigl(F(X_{2})^{\vee},\ F(X_{1})\bigr)
\]
Familles : $X_{I} = \prod_{j} X_{I_{j}}$, $I = I^{+} \amalg I^{-}$,
\[
F(X_{I})
\simeq \underline{\mathrm{Hom}}^{(1)}\Bigl(\prod_{j\in J} F(X_{I_{j}})^{\circ},\
\mathrm{Ens}\Bigr)
\simeq \underline{\mathrm{Hom}}_{(1)}\Bigl(\prod_{j\in J}
F(X_{I_{j}}^{\vee}),\ \mathrm{Ens}\Bigr)
\simeq \cdots
\]
\[
\cdots \simeq \underline{\mathrm{Hom}}_{!}\Bigl(\prod_{j\in J^{-}}
F(X_{I_{j}}),\ \prod_{j'\in J^{+}} F(X_{I_{j'}})\Bigr)
\simeq \underline{\mathrm{Hom}}^{!}\Bigl(\prod_{j'\in J^{+}}
F(X_{I_{j'}})^{\circ},\ \prod_{j\in J^{-}} F(X_{I_{j}}^{\vee})\Bigr)
\]
\note{décorations ① et exposants en partie conjecturaux}
En bas : $\mathrm{Top}(A) \xrightarrow{\varphi} X \xleftarrow{\Psi}
\mathrm{Top}(B)$, \quad $A \to \mathrm{Pt}(X)$
\note{trois lignes finales raturées, mêlant $\varphi_{*}(F)$,
$F \in \mathrm{Top}(A)$, $G \in \mathrm{Ob}\,B^{\vee}$ — lecture trop
incertaine pour être poursuivie}

\page{9}
\note{page 4 de la main de l'auteur ; la page ne porte que ces deux
formules, encadrées}
\[
\bigl\langle \Phi_{*},\ f^{\circ*}(\Psi^{*}) \bigr\rangle_{A}
\simeq \bigl\langle f_{!}(\Phi_{*}),\ \Psi^{*} \bigr\rangle_{B},
\qquad
\bigl\langle f^{*}(\Psi_{*}),\ \Phi^{*} \bigr\rangle_{A}
\simeq \bigl\langle \Psi_{*},\ f^{\circ}_{!}(\Phi^{*}) \bigr\rangle_{B}
\]
\note{répartition des étoiles et des ronds partiellement conjecturale}

\section*{L'enveloppe $\widetilde{A}$ (pages 10 et 11)}

\page{10}
\note{page 5 de la main de l'auteur}
« $A \hookrightarrow B$ \uncertain{pl.\ fid.}, $x \in B$ :
\[
H_{x} \in \widehat{A}, \quad H_{x}(a) = \mathrm{Hom}_{B}(a, x) ;
\qquad
H^{x} \in A^{\vee}, \quad H^{x}(a) = \mathrm{Hom}_{B}(x, a)
\]
\[
H_{x}(a) \times H^{x}(b) \xrightarrow{\ \alpha_{x}\ } \mathrm{Hom}_{A}(a, b)
\]
 »
\note{la flèche est la composition dans $B$ ; elle aboutit dans
$\mathrm{Hom}_{A}$ parce que $A$ est pleine dans $B$}

« Soit $\widetilde{A}$ la catégorie des triples $(H_{*}, H^{*}, \alpha)$ où
$H_{*} \in \widehat{A}$, $H^{*} \in A^{\vee}$,
$\alpha : H_{*}(a) \times H^{*}(b) \to \mathrm{Hom}_{A}(a,b)$ (fonctoriel).
Les flèches $(H_{*}, H^{*}, \alpha) \xrightarrow{\varphi}
(H'_{*}, H'^{*}, \alpha')$ étant \uncertain{formées des couples}
$\varphi_{*} : H_{*} \to H'_{*}$ dans $\widehat{A}$,
$\varphi^{*} : H'^{*} \to H^{*}$ dans $A^{\vee}$, tels que l'on ait :
\[
\alpha_{a,b}\bigl(u,\ \varphi^{*}(u')\bigr)
= \alpha'_{a,b}\bigl(\varphi_{*}(u),\ u'\bigr)
\]
 »
\note{le manuscrit exprime cette identité par un carré à variance mixte,
$\varphi_{*}$ descendant sur le premier facteur et $\varphi^{*}$ remontant
sur le second}

« Pour un foncteur \uncertain{pl.\ fid.} $A \xrightarrow{\varphi} B$, on en
\uncertain{déduit} un foncteur $\varphi^{\S} : B \to \widetilde{A}$,
$x \mapsto (\varphi^{\S}_{*}(x),\ \varphi^{\S *}(x),\ \alpha_{x})$. »
\note{décorations des deux composantes incertaines}

\page{11}
\note{page 6 de la main de l'auteur}
« NB $\widetilde{A}^{\circ} \simeq \widetilde{A}$ \ldots »
\note{fin de ligne illisible}
« Cas où $\varphi^{\S}$ est pl.\ fidèle ? Il suffit que $f$ soit
\uncertain{strict.\ gén.} ou \ill{} \uncertain{adjonctions}. »

« Cas $A \xrightarrow{\mathrm{can}} \widehat{A}$ :
$\varphi^{\S}(F)_{*} = F$,\quad
$\varphi^{\S}(F)^{*}(a) = \mathrm{Hom}(F, a)$ ;\quad $\varphi^{\S}$
\uncertain{fidèle}. »

Petits diagrammes intermédiaires sur $F$, $F^{\S}$, $G$, avec $u \in F(a)$,
$u' \in F^{\S}(b)$.
\note{enchaînement raturé ; on y lit « $v = v' \circ \varphi_{*}$ ! » et
« on sait que $v \circ u = v' \circ u'$ »}

\struck{Considérons alors la sous-catégorie pleine de $\widetilde{A}$ formée
des} \ill{}
\note{passage biffé puis récrit, en grande partie illisible ; y figurent les
chaînes $\Psi_{1}^{\circ} \to a \to \Phi \to \Psi_{2}^{\circ}$ et
$\Phi_{1} \to \Psi \to \Phi_{2}$}

$(F, G, \alpha) \simeq (F, F^{\S}, \alpha_{F})$, \qquad
$(F, G, \alpha) \simeq (G^{\S}, G, \alpha^{G})$
\note{lecture conjecturale des deux isomorphismes}

Fonctorialité : $A \xrightarrow{f} B$ induit
$\widetilde{f} : \widetilde{A} \to \widetilde{B}$ ; en bas de la page,
encerclé :
\begin{tikzcd}[column sep=normal, row sep=normal]
\widehat{A} \arrow[r, hook] \arrow[d, "f_{!}"']
  & \widetilde{A} \arrow[d, Rightarrow, "\widetilde{f}"]
  & A^{\vee\circ} \arrow[l, hook'] \arrow[d, "f^{\circ*\circ}"] \\
\widehat{B} \arrow[r, hook]
  & \widetilde{B}
  & B^{\vee\circ} \arrow[l, hook']
\end{tikzcd}
\note{étiquettes des flèches verticales partiellement conjecturales ; à côté,
une colonne de paires de foncteurs adjoints — $f_{!}, f^{*}$ ;
$f^{\circ}_{!}, f^{\circ *}$ ; \ldots{} — trop abrégée pour être restituée}

\section*{Noyau autodual, recollements, coefficients (pages 13 et 14)}

\page{13}
\note{page 7 de la main de l'auteur}
En haut :
\[
\widehat{A} \xrightarrow{\ \approx\ }
\underline{\mathrm{Hom}}_{!}(A^{\vee}, \mathrm{Ens}),
\qquad
A^{\vee\circ} \xrightarrow{\ \approx\ }
\underline{\mathrm{Hom}}_{\mathrm{rep}}(\widehat{A}, (\mathrm{Ens}))
\]
\note{l'indice du second Hom — « rep » ? — est incertain ; au-dessus des deux,
un $\underline{\mathrm{Hom}}(A^{\vee}, (\mathrm{Ens}))$ sans décoration les
reçoit tous deux}
Au centre, relié à $A$ par une double flèche verticale :
\[
\bigl(\mathrm{Kar}(A^{\circ})\bigr)^{\circ} \;\simeq\; \mathrm{Kar}(A)
\]
À droite, le triangle $\varepsilon(A)$ sur $\widehat{A}$ et $A^{\vee\circ}$,
avec
\[
A = \widehat{A} \cap A^{\vee\circ}
\]
\note{à gauche de la page, un $\overline{\varepsilon}(A)$ isolé, sous un
trait de séparation}

En bas, la mise en garde sur les recollements :
« $N = M_{1} \cap M_{2}$ \uncertain{via} des flèches pl.\ fid. Alors on
\uncertain{définit} une catégorie $\Sigma'$, \uncertain{réunion} des
sous-catégories pleines $M'_{1}$ et $M'_{2}$, avec \uncertain{intersection}
$N'$, et un foncteur pl.\ fid.\ $\Sigma' \to M$ (\uncertain{factorisant par}
$M_{1} \cup M_{2}$). Mais \uncertain{même} si $M_{1}$, $M_{2}$ sont pleines
dans $M$ (avec $N$ \ill{} dans $M$) et $M'_{1} \to M_{1}$,
$M'_{2} \to M_{2}$ des équivalences, $M'_{1}$ et $M'_{2}$ ne sont \ill{}
\ill{} pleines dans $\Sigma'$, et $N' \to N$ n'est pas \ill{} une
équivalence. »
\note{la lecture de plusieurs mots reste incertaine}

\page{14}
Dualité à valeurs dans une catégorie $M$ :
\[
\underline{\mathrm{Hom}}(A, M)
\simeq \underline{\mathrm{Hom}}_{!}(\widehat{A}, M)
\]
\note{un troisième membre, aux décorations illisibles, prolonge la ligne}
« Ex.\ $M = $ \ill{} : $\underline{\mathrm{Hom}}(B^{\circ}, M) \simeq \cdots$ »
\note{exemple raturé}

« $A = B^{\circ}$, $B = A^{\circ}$ » :
\[
\underline{\mathrm{Hom}}(A, M) \simeq
\underline{\mathrm{Hom}}^{!}(\widehat{A}, M)
\]
« \uncertain{co}faisceaux sur $A$, ind.\ dans $M$ » ;
\[
\underline{\mathrm{Hom}}(B^{\circ}, M) \simeq
\underline{\mathrm{Hom}}^{!}(\widehat{B}^{\circ}, M)
\]
« faisceaux sur $B$, ind.\ dans $M$ ».
\[
\mathcal{F}(X, M)^{\circ} \simeq \check{\mathcal{F}}(X, M^{\circ}),
\qquad
\bigl(\underline{\mathrm{Hom}}^{!}(A, M)\bigr)^{\circ}
\simeq \underline{\mathrm{Hom}}_{!}(A^{\circ}, M^{\circ})
\]
« $\mathcal{F}(X,M)$ et $\check{\mathcal{F}}(X, M')$ \uncertain{se déduisent
l'un de l'autre} » ;
\[
\mathcal{F}(X, M) \approx \check{\mathcal{F}}(X^{\circ}, M),
\qquad
\underline{\mathrm{Hom}}^{!}(A^{\circ}, M)
\simeq \underline{\mathrm{Hom}}_{!}(A, M^{\circ})^{\circ}
\]
\note{dernières lignes aux décorations incertaines}

\end{document}
